% فصل ۲ — مرور ادبیات
% همهٔ عددها از docs/generated/thesis_numbers.json می‌آیند و در results/VERIFICATION-*.md
% در برابر منبع اصلی راستی‌آزمایی شده‌اند. ادعای راستی‌آزمایی‌نشده در این فصل نقل نمی‌شود.

\chapter{مرور ادبیات}
\label{ch:related}

این فصل ادبیات را بر پایهٔ موضوع سازمان نمی‌دهد، بلکه بر پایهٔ \textbf{تناقض} سازمان می‌دهد.
دلیلش این است که فهرستی از «چه کسی چه کرد» به پرسش‌های فصل \ref{ch:intro} پاسخ نمی‌دهد، در
حالی که جاهایی که نتیجه‌های منتشرشده با هم نمی‌خوانند دقیقاً همان‌جاهایی‌اند که یک آزمایش تازه
چیزی به دست می‌آورد.

سه تناقض بررسی می‌شود. در هر سه، الگو یکسان است: \textbf{نتیجه‌های ناسازگار در خانه‌هایی
نشسته‌اند که مقالهٔ مبدأ آن‌ها را نیازموده است.} پس تناقض‌ها با انتخاب یک طرف حل نمی‌شوند، با
نام بردن خانهٔ آزموده‌نشده حل می‌شوند.

\paragraph{یادداشتی دربارهٔ اتکاپذیری این فصل.} هر ادعای نقل‌شده در این فصل در برابر متن اصلی
مقاله بررسی شده است و پرونده‌های راستی‌آزمایی در مخزن پژوهش نگهداری می‌شوند. جایی که عددی در
منبع پیدا نشد یا با آنچه انتظار می‌رفت نخواند، \textbf{آن عدد در این متن نقل نمی‌شود}. سه مورد
از این دست وجود دارد و در بخش \ref{sec:lit-corrections} نام برده می‌شوند، چون خودِ آن‌ها
یافته‌اند.

\section{مدلی که زیر حسابرسی است}
\label{sec:retfound-audit}

نقطهٔ مرجع این ادبیات یک مدل بنیادین شبکیه است که با یادگیری خودنظارتی روی \textbf{\num{۱٫۶}}
میلیون تصویر ته‌چشم و مقطع‌نگاری همدوسی نوری پیش‌آموزش دیده و برای وظایف پایین‌دستی ریزتنظیم
می‌شود \cite{zhou2023retfound}. مروری داوری‌شده آن را چنین خلاصه می‌کند: این مدل خودنظارتی
مدل‌های نظارت‌شده و خودنظارتیِ پیش‌آموزش‌دیده روی \lr{ImageNet} را در همان وظایف، \emph{هم در
اعتبارسنجی درونی و هم بیرونی}، پشت سر می‌گذارد \cite{ruamviboonsuk2024review}.

\textbf{این ادعا زیر حسابرسی است، نه زیر سؤال.} تفاوت مهم است. مقالهٔ مبدأ آنچه را آزموده
درست گزارش می‌کند؛ آنچه این فصل نشان می‌دهد این است که \emph{بخش بزرگی از فضای مسئله را
نیازموده}، و هر نتیجهٔ ناسازگارِ بعدی دقیقاً در همان بخش نشسته است. مشخصاً: در آن مقاله هیچ
پایهٔ شبکهٔ همگشتی وجود ندارد، هیچ رمزگذار منجمد یا کاوشگر خطی ارزیابی نمی‌شود، هیچ معیار
به‌تفکیک‌کلاس برای رتینوپاتی گزارش نمی‌شود، و واسنجی زیر جابه‌جایی توزیع سنجیده نمی‌شود.

\section{تناقض یکم: مدل بنیادین منجمد در برابر ریزتنظیم‌شده}
\label{sec:contradiction1}

\subsection{شواهدی که با هم نمی‌خوانند}

سه گروه مستقل، سه پاسخ متفاوت به «آیا مدل بنیادین شبکیه بهتر است؟» می‌دهند.

\textbf{پاسخ یکم: بله، و تفاوت ناچیز است.} روی یک پیکرهٔ چندپیمانه‌ای، دقت درجه‌بندی
رتینوپاتی روی تصویر ته‌چشم برای مدل بنیادین \textbf{\num{۰٫۸۲۲}} و برای یک شبکهٔ همگشتی
پیش‌آموزش‌دیده روی \lr{ImageNet} برابر \textbf{\num{۰٫۸۲۳}} است — عملاً مساوی
\cite{tang2026mmrdr}.

\textbf{پاسخ دوم: بله، ولی فقط درون توزیع.} یک سر لجستیک با \textbf{\num{۵۱۲۵}} پارامتر روی
ویژگی‌های منجمد همان مدل بنیادین، با ریزتنظیم کامل همان مدل که \textbf{\num{۳۰۳٫۳}} میلیون
پارامتر دارد تفکیک‌ناپذیر است \cite{nielsen2025he}. اعداد این‌ها هستند: روی
\lr{APTOS} سر منجمد \textbf{\num{۰٫۹۳}} در برابر ریزتنظیم \textbf{\num{۰٫۹۴}}، و روی \lr{ODIR-5K}
\textbf{\num{۰٫۷۸}} در برابر \textbf{\num{۰٫۸۰}}. \textbf{در هر دو مورد سر منجمد اسماً عقب‌تر است و
اختلاف معنادار نیست.} نویسندگان خود می‌نویسند که ارزیابی بیرون از توزیع را انجام نداده‌اند و
انجامش بینش ارزشمندی می‌داد.

\textbf{پاسخ سوم: نه، وقتی پیمانه عوض شود.} روی همان پیکرهٔ چندپیمانه‌ای، هنگام انتقال به
تصویربرداری فوق‌عریض، شبکهٔ همگشتی پیش‌آموزش‌دیده روی \lr{ImageNet} \emph{بهتر} از مدل‌های
بنیادین چشم‌پزشکی تعمیم می‌دهد \cite{tang2026mmrdr}.

\subsection{چرا این‌ها در واقع تناقض نیستند}

هر سه نتیجه در خانه‌ای نشسته‌اند که مقالهٔ مبدأ نیازموده است. مقالهٔ مبدأ رمزگذار منجمد را
ارزیابی نمی‌کند، پس با پاسخ دوم برخورد نمی‌کند. پایهٔ همگشتی ندارد، پس با پاسخ‌های یکم و سوم
برخورد نمی‌کند. \textbf{تناقض ظاهری، محصول تعمیم دادن نتیجه‌ای است به شرایطی که در آن آزموده
نشده.}

نتیجهٔ خود این پژوهش در فصل \ref{ch:results} در همین شکاف می‌نشیند و آن را تیزتر می‌کند: با
رمزگذار منجمد، بازنمایی مدل بنیادین \emph{بهتر} است؛ با ریزتنظیم کامل، همان مدل به خط لولهٔ
ساده‌تر ما \emph{می‌بازد}. دو سازوکار نامزد وجود دارد و شواهد ما آن‌ها را از هم جدا نمی‌کند؛
این در فصل \ref{ch:discussion} صریح گفته می‌شود.

\subsection{یک اشکال روش‌شناختی که در همین‌جا تکرار می‌شود}

پیکرهٔ چندپیمانه‌ای یادشده مدل‌ها را کنار هم جدول می‌کند در حالی که \emph{شیوهٔ تطبیق آن‌ها
یکسان نیست}: برخی به‌طور کامل ریزتنظیم می‌شوند و برخی با «تنظیم حداقلی» ارزیابی می‌شوند. این
دقیقاً همان مخدوش‌کننده‌ای است که فصل \ref{ch:method} در بخش \ref{sec:matched} می‌بندد.

\textbf{و نمونهٔ خلافش در ادبیات وجود دارد.} مقالهٔ مبدأ، هنگام مقایسهٔ راهبردهای پیش‌آموزش،
معماری و فرایند تطبیق را ثابت نگه می‌دارد و تنها چیز تحت آزمون را تغییر می‌دهد، و خود هشدار
می‌دهد که ادعای برتری یک روش با وجود چند متغیر هم‌زمان نیازمند احتیاط است
\cite{zhou2023retfound}. \textbf{این طراحی درست است و باید معیار سنجش بقیه باشد.}

\section{تناقض دوم: تشخیص‌پذیری به‌جای واسنجی}
\label{sec:contradiction2}

\subsection{خطا در طبیعت}

یک مقالهٔ داوری‌شدهٔ \num{۲۰۲۶}، پس از گزارش مقادیر \lr{AUC}، می‌نویسد که این مقادیر بالا «تأیید
می‌کنند» مدل نه‌تنها دقیق است بلکه احتمال‌های خوش‌واسنجی نسبت می‌دهد و برای پشتیبانی تصمیم
بالینی قابل اتکاست \cite{yu2026dualswinord}.

\textbf{این گزاره نادرست است، و نادرستی‌اش ساختاری است نه کمّی.} \lr{AUC} تنها به
\emph{ترتیب} نمره‌ها حساس است و نسبت به هر تبدیل یکنوا ناوردا است. مدلی که همهٔ احتمال‌هایش را
در عدد ثابتی ضرب کند، \lr{AUC} یکسانی دارد و واسنجی کاملاً متفاوتی. آن مقاله \emph{هیچ} معیار
واسنجی گزارش نمی‌کند — نه \lr{ECE}، نه \lr{Brier}، نه منحنی اطمینان — و هیچ اعتبارسنجی بیرونی
هم ندارد: هر دو محکش، پیکره‌هایی‌اند که روی آن‌ها آموزش دیده است.

\subsection{شاهد نقض، از دل همان ادبیات}

مسئله را نمی‌توان با شهود رد کرد؛ با اندازه‌گیری رد می‌شود، و اندازه‌گیری‌اش منتشر شده است. یک
محک واسنجی‌محور نشان می‌دهد که همان مدل بنیادین، \emph{کمترین} خطای واسنجی بیرونی را در کنار
\emph{بدترین} تشخیص‌پذیری بیرونی به دست می‌آورد \cite{poyrazer2026calibration}. دو معیار در
جهت مخالف هم حرکت می‌کنند، پس هیچ‌کدام جانشین دیگری نیست.

\textbf{و واسنجی زیر جابه‌جایی توزیع فرو می‌ریزد.} در یک محک ترتیبی، خطای واسنجی از
\textbf{\num{۰٫۰۴۹}} روی مجموعهٔ درونی به \textbf{\num{۰٫۱۶۰}} روی مجموعهٔ بیرونی می‌رسد
\cite{sheng2026orderdr}. سه برابر شدن خطای واسنجی، بدون آن‌که مدل عوض شود.

\section{تناقض سوم: ادعاهای معماری که برمی‌گردند}
\label{sec:contradiction3}

یک مرور داوری‌شده می‌نویسد پژوهش‌های بسیاری برتری ترانسفورمر بینایی را بر شبکهٔ همگشتی برای
وظایف رایج، از جمله غربالگری رتینوپاتی، یافته‌اند. \emph{در همان مرور} و بدون آشتی دادن این
دو، مقایسه‌ای سربه‌سر میان هشت مدل همگشتی و نه مدل ترانسفورمر گزارش می‌شود که در آن
\textbf{همهٔ} مدل‌های همگشتی از \textbf{همهٔ} مدل‌های ترانسفورمر بهترند: حساسیت و ویژگی
مدل‌های همگشتی به \num{۹۰} درصد یا بیشتر می‌رسد، در حالی که ترانسفورمرها به حساسیت \num{۶۳} تا \num{۹۴} درصد و
ویژگی \num{۲۴} تا \num{۴۸} درصد می‌رسند \cite{ruamviboonsuk2024review}.

\textbf{این تناقض در همان صفحه است و آشتی داده نمی‌شود.} خواندن درست آن این نیست که یکی از دو
گزاره غلط است؛ این است که \textbf{ادعای معماری، وقتی شرایط ارزیابی عوض شود، برمی‌گردد}. نتیجهٔ
خود این پژوهش در فصل \ref{ch:results} همین را به شکل کمّی می‌گوید: از نه مداخله، هشت مداخله
ابطال شدند، و در یک مورد \emph{علامت} اختلاف با جابه‌جایی نقاط برش وارونه شد.

\section{خوشهٔ دوم: شواهد استقرار}
\label{sec:deployment}

خوشهٔ دوم ادبیات، مقالهٔ توسعهٔ مدل نیست؛ شواهد میدانی است. اهمیتش این است که استدلال قاعدهٔ
تصمیم را از یک نکتهٔ روش‌شناختی به یک مسئلهٔ بالینی تبدیل می‌کند.

\textbf{نقطهٔ عمل، نه مدل، خروجی را تعیین می‌کند.} در اعتبارسنجی واقعی سه سامانهٔ تجاری روی
همان چشم‌ها و با همان مرجع، ویژگی از \textbf{\num{۱۴٫۲۵}} تا \textbf{\num{۹۶٫۰۱}} درصد نوسان می‌کند
\cite{duggal2025realworld}. در همان کار، شاخص توافق از \textbf{\num{۰٫۶۵}} به \textbf{\num{۰٫۷۲}}
می‌رود در حالی که حساسیت و ویژگی هر کدام حدود سی‌ویک واحد جابه‌جا می‌شوند.

\paragraph{و این جابه‌جایی را نباید جاروب آستانه خواند.} دو مقدار توافق از دو فاز متفاوت، دو
گروه بیمار متفاوت و پس از بازآموزی الگوریتم توسط فروشنده می‌آیند — خودِ مقاله این را می‌گوید.
\textbf{پس آن را به‌عنوان شاهدِ بی‌حسی آماره‌های توافق تجمیعی نسبت به تغییر نقطهٔ عمل نقل
می‌کنیم، نه به‌عنوان جاروب آستانه روی دادهٔ ثابت.} این تمایز را در جای دیگری از این متن نباید
از دست داد.

\textbf{مرز خطا جای مشخصی دارد.} در پژوهشی روی \textbf{\num{۳۲۰}} چشم از \textbf{\num{۱۶۰}} بیمار، سه
روش غربالگری در تشخیص درجه‌های متوسط و شدید و تکثیری هیچ تفاوتی ندارند، و تمام اختلافشان در
چشم‌های بدون رتینوپاتی و چشم‌های با نشانه‌های خفیف است \cite{tomic2024handheld}. توافق
\textbf{\num{۰٫۸۶}} در برابر معاینهٔ بالینی و \textbf{\num{۰٫۷۸}} در برابر دوربین استاندارد است، و
اختلاف‌ها \textbf{\num{۲۲}} و \textbf{\num{۳۴}} چشم‌اند.

\textbf{و سر ادم ماکولا در بالین همان‌جایی می‌شکند که در آزمایشگاه.} یک سامانهٔ تجاری مستقر،
در یک اجرا و روی همان تصویرها، برای رتینوپاتی نیازمند ارجاع به توافق \textbf{\num{۰٫۸۱}} می‌رسد و
برای ادم ماکولا روی تصویر ته‌چشم به حساسیت \textbf{\num{۲۶٫۵}} درصد و توافق \textbf{\num{۰٫۳۸}}
\cite{duggal2025realworld}. \textbf{این تک‌شاهد، همهٔ توضیح‌های جایگزین برای سقف ادم ماکولا را
می‌بندد}: نه معماری ما، نه اندازهٔ دادهٔ ما، نه بودجهٔ آموزش ما، نه انتخاب پیکره.

\section{شکاف تبار داده}
\label{sec:lit-provenance}

آخرین موضوع این فصل، موضوعی است که ادبیات به‌ندرت دربارهٔ خودش گزارش می‌کند.

یک فراتحلیل، حساسیت تجمیعی \textbf{\num{۰٫۹۵}} گزارش می‌کند در حالی که پژوهش‌های واردشده بازه‌ای
تا \textbf{\num{۰٫۰۶}} را در بر می‌گیرند \cite{tahir2025meta}؛ و بررسی جدول‌های همان کار نشان
می‌دهد ستون منفی‌های کاذب رونویسی ستون مثبت‌های کاذب است، و همان خرابی به حساسیت‌های
گزارش‌شده نشت کرده است. یک توصیف‌نامهٔ داده در خانوادهٔ \lr{Nature} پیکرهٔ خود را از یک
مجموعهٔ عمومی می‌سازد که در پیش‌آموزش یکی از مدل‌های محک‌شده به کار رفته، بی‌آن‌که این
هم‌پوشانی را جایی بیاورد \cite{tang2026mmrdr}؛ و زیرمجموعه‌های خود را در دو پیمانه در سطح
بیمار تقسیم می‌کند «برای جلوگیری از نشت داده» و در سومی در سطح تصویر، چون پیکرهٔ مبدأ شناسهٔ
بیمار ندارد. یک پژوهش بالینی، \textbf{\num{۳۲۰}} چشم از \textbf{\num{۱۶۰}} بیمار را چشم‌به‌چشم تحلیل
می‌کند بی‌آن‌که برای همبستگی دو چشم یک نفر تصحیحی بگذارد \cite{tomic2024handheld}.

\textbf{هیچ‌کدام از این‌ها خطای زمان اجرا تولید نمی‌کند.} همه در سکوت عدد می‌سازند. فصل
\ref{ch:method} سازوکارهایی را می‌آورد که در این پژوهش پاسخ آن‌هاست، و در بخش
\ref{sec:our-own-failure} \textbf{شکست هم‌جنس خودِ این پایان‌نامه} را در کنار همین چهار نمونه
می‌گذارد.

\section{سه تصحیح که خودشان یافته‌اند}
\label{sec:lit-corrections}

راستی‌آزمایی این فصل در برابر متن اصلی، سه موردی را پیدا کرد که نقلشان بدون تصحیح، خطا را
تکثیر می‌کرد. ثبتشان اینجا آمده چون \emph{نمونهٔ همان مسئله‌ای‌اند که این فصل دربارهٔ آن است}.

\textbf{یکم، معیار عوض می‌شود.} مقادیر \textbf{\num{۰٫۸۲۲}} و \textbf{\num{۰٫۸۲۳}} در بخش
\ref{sec:contradiction1} \emph{دقت}‌اند، نه \lr{AUC}. آن مقاله صریحاً می‌نویسد که \lr{AUC} را
به دلیل هزینهٔ محاسباتی کنار گذاشته است. نقل کردنشان به‌عنوان \lr{AUC} معیاری را به مقاله
نسبت می‌داد که خودش از محاسبه‌اش سر باز زده است.

\textbf{دوم، جهت عوض می‌شود.} در مقایسهٔ سر منجمد با ریزتنظیم، سر منجمد روی هر دو مجموعه
اسماً \emph{عقب‌تر} است، نه هم‌تراز یا جلوتر. نتیجهٔ «تفکیک‌ناپذیر» سرجایش می‌ماند، ولی جهتی
که خلاصه‌های ثانویه القا می‌کنند در منبع وجود ندارد.

\textbf{سوم، نرخی که بر پایهٔ سال اشتباه بنا شده.} مرور یادشده پس از گزارش درست
\textbf{\num{۳}} دستگاه دارای مجوز، نرخ «یک مدل در سال» را از این فرض می‌گیرد که نخستین مجوز در
\num{۲۰۲۱} صادر شده است؛ نخستین مجوز در \num{۲۰۱۸} صادر شد. \textbf{شمار دستگاه‌ها نقل می‌شود، نرخ نه.}

یک ادعای چهارم هم هست که در این پایان‌نامه \emph{اصلاً} نقل نمی‌شود، چون در متن اصلی یافت نشد
و متن اصلی ظاهراً خلافش را می‌گوید. تصمیم بر این شد که به‌جای بازنویسی محتاطانه، کنار گذاشته
شود: \textbf{ادعایی که راستی‌آزمایی نشده، نقل نمی‌شود.}

\section{جای این پژوهش}

سه تناقض این فصل، سه خانهٔ آزموده‌نشده را نام می‌برند: مقایسهٔ منجمد و ریزتنظیم‌شده در یک
طراحی که تنها یک چیز را عوض کند؛ واسنجی به‌عنوان معیاری مستقل از تشخیص‌پذیری و زیر جابه‌جایی
توزیع؛ و ادعای معماری زیر نقاط برش همتاشده. فصل \ref{ch:method} پروتکلی را می‌سازد که هر سه
را هم‌زمان می‌بندد، و فصل \ref{ch:results} نشان می‌دهد وقتی این کار انجام شود، بیشتر
مداخله‌ها چیزی نمی‌سازند — و همین، نتیجه است.
